% Rapporto esecutivo: compilare con xelatex (due passaggi).
% Prima della consegna: sostituire ogni segnaposto <<...>> e cancellare le righe % GUIDA.
% Ordine delle sezioni: che cosa, perché, come, che cosa abbiamo trovato,
% interpretazione e limiti, note di versione (solo edizioni riviste).

% Configurazione
\newcommand{\ReportLang}{italian}% italian | english
\newcommand{\ReportBrand}{Nomisma}
\newcommand{\ReportProject}{<<PROGETTO>>}% testatina e piede; vuoto per il solo marchio
\newcommand{\ReportTitle}{<<TITOLO: il risultato in una riga, con la scala>>}
\newcommand{\ReportSubtitle}{<<SOTTOTITOLO: Rapporto di sintesi sul lavoro svolto e sui risultati>>}
\newcommand{\ReportAuthor}{<<NOME COGNOME>>}
\newcommand{\ReportRole}{<<RUOLO>>}
\newcommand{\ReportDate}{<<GIORNO MESE ANNO>>}

% report-base.tex: preambolo comune per i rapporti PDF in stile nota Nomisma.
% Uso: nel file principale definire le macro di configurazione (\ReportLang,
% \ReportBrand, \ReportProject, \ReportTitle, \ReportSubtitle, \ReportAuthor,
% \ReportRole, \ReportDate), poi % report-base.tex: preambolo comune per i rapporti PDF in stile nota Nomisma.
% Uso: nel file principale definire le macro di configurazione (\ReportLang,
% \ReportBrand, \ReportProject, \ReportTitle, \ReportSubtitle, \ReportAuthor,
% \ReportRole, \ReportDate), poi % report-base.tex: preambolo comune per i rapporti PDF in stile nota Nomisma.
% Uso: nel file principale definire le macro di configurazione (\ReportLang,
% \ReportBrand, \ReportProject, \ReportTitle, \ReportSubtitle, \ReportAuthor,
% \ReportRole, \ReportDate), poi \input{report-base} e \begin{document}.
% Compilare con xelatex, due passaggi (tre se cambiano riferimenti o bibliografia).
\documentclass[10pt,a4paper]{article}

% Valori di default: il file principale li sovrascrive definendoli prima di \input.
\providecommand{\ReportLang}{italian}% italian | english
\providecommand{\ReportBrand}{Nomisma}
\providecommand{\ReportProject}{}
\providecommand{\ReportTitle}{<<TITOLO>>}
\providecommand{\ReportSubtitle}{<<SOTTOTITOLO>>}
\providecommand{\ReportAuthor}{<<AUTORE>>}
\providecommand{\ReportRole}{<<RUOLO>>}
\providecommand{\ReportDate}{<<DATA>>}

\usepackage{etoolbox}
\usepackage{fontspec}
% Testo estraibile corretto dal PDF: senza questa riga le parentesi di Segoe UI
% vengono copiate come caratteri ornamentali (U+FD3E/U+FD3F).
\XeTeXgenerateactualtext=1

% Font: Segoe UI di Windows come nella nota di riferimento; fuori da Windows si
% ripiega su Latin Modern Sans, sempre presente in MiKTeX e TeX Live.
\IfFileExists{C:/Windows/Fonts/segoeui.ttf}{%
  \setmainfont{segoeui}[
    Path = C:/Windows/Fonts/,
    Extension = .ttf,
    BoldFont = segoeuib,
    ItalicFont = segoeuii,
    BoldItalicFont = segoeuiz]
  \setsansfont{segoeui}[
    Path = C:/Windows/Fonts/,
    Extension = .ttf,
    BoldFont = segoeuib,
    ItalicFont = segoeuii,
    BoldItalicFont = segoeuiz]
}{%
  \setmainfont{Latin Modern Sans}
  \setsansfont{Latin Modern Sans}
}

\usepackage[english,italian]{babel}
\AtBeginDocument{\selectlanguage{\ReportLang}}

\usepackage[a4paper,top=2.3cm,bottom=2.2cm,left=2.2cm,right=2.2cm,headheight=14pt,headsep=0.6cm,footskip=1cm]{geometry}
\usepackage[table]{xcolor}
\usepackage{colortbl}
\usepackage{array}
\usepackage{tabularx}
\usepackage{xltabular}
\usepackage{titlesec}
\usepackage{fancyhdr}
\usepackage{enumitem}
\usepackage{microtype}
\usepackage{needspace}
\usepackage{titletoc}
\usepackage[section]{placeins}
\usepackage{amsmath}
\usepackage{amssymb}
\usepackage{graphicx}
\usepackage{caption}
\usepackage{siunitx}

% Colori della nota metodologica
\definecolor{navy}{HTML}{0E2D53}
\definecolor{maroon}{HTML}{A3274B}
\definecolor{rowgrey}{HTML}{F2F2F2}
\definecolor{rulegrey}{HTML}{BFBFBF}
\definecolor{footgrey}{HTML}{8C8C8C}
\definecolor{boxgrey}{HTML}{F4F5F8}

\setlength{\parindent}{0pt}
\setlength{\parskip}{0.55em}
\linespread{1.12}

% Numeri: in italiano punto per le migliaia e virgola per i decimali, in inglese
% il contrario. Nel sorgente si scrive sempre il numero grezzo: \num{161250},
% \num{0.698}, \pct{96.6}. Il raggruppamento parte da 4 cifre (1.626).
% Modo testo: le cifre usano il font del testo e ereditano il grassetto.
% Incertezza esplicita: \num{0.874 +- 0.006} produce 0,874 ± 0,006.
\sisetup{mode=text,reset-text-family=false,reset-text-series=false,reset-text-shape=false,uncertainty-mode=separate}
\ifdefstring{\ReportLang}{english}{%
  \sisetup{output-decimal-marker={.}, group-separator={,}, group-minimum-digits=4}
  \newcommand{\ReportPageLabel}{p.}
}{%
  \sisetup{output-decimal-marker={,}, group-separator={.}, group-minimum-digits=4}
  \newcommand{\ReportPageLabel}{pag.}
}
\newcommand{\pct}[1]{\num{#1}\%}

% Testatina bordeaux in maiuscolo a sinistra, piede grigio a destra
\providecommand{\ReportHeader}{\ReportBrand\ifdefempty{\ReportProject}{}{\ \ | \ \ReportProject}}
\providecommand{\ReportFooter}{\ReportBrand\ifdefempty{\ReportProject}{}{ | \ReportProject}}
\newcommand{\ReportHeadStyle}{\footnotesize\bfseries\color{maroon}\MakeUppercase{\ReportHeader}}
\newcommand{\ReportFootStyle}{\scriptsize\color{footgrey}\ReportFooter\ ·\ \ReportPageLabel\ \thepage}
\pagestyle{fancy}
\fancyhf{}
\renewcommand{\headrulewidth}{0pt}
\fancyhead[L]{\ReportHeadStyle}
\fancyfoot[R]{\ReportFootStyle}
\fancypagestyle{plain}{\fancyhf{}\renewcommand{\headrulewidth}{0pt}%
  \fancyhead[L]{\ReportHeadStyle}\fancyfoot[R]{\ReportFootStyle}}

% Blocco del titolo: titolo blu in grassetto, sottotitolo, autore e ruolo, data
\newcommand{\maketitleblock}{%
  {\color{navy}\bfseries\LARGE \ReportTitle\par}
  \vspace{0.35em}
  {\color{navy}\large \ReportSubtitle\par}
  \vspace{0.5em}
  {\small \ReportAuthor, \ReportRole\par
  \color{footgrey}\ReportDate\par}}

% Titoli di sezione: numero e titolo blu in grassetto con filetto bordeaux sotto.
% Le sottosezioni sono blu, in grassetto, di corpo normale (piu' piccole delle sezioni).
\newcommand{\sectitle}[1]{\parbox[b]{\linewidth}{\thesection.\hspace{0.4em}#1\par\nointerlineskip\vspace{3pt}{\color{maroon}\rule{\linewidth}{0.7pt}}}}
\newcommand{\sectitlenn}[1]{\parbox[b]{\linewidth}{#1\par\nointerlineskip\vspace{3pt}{\color{maroon}\rule{\linewidth}{0.7pt}}}}
\titleformat{\section}[block]{\color{navy}\bfseries\large}{}{0pt}{\sectitle}
\titleformat{name=\section,numberless}[block]{\color{navy}\bfseries\large}{}{0pt}{\sectitlenn}
\titlespacing*{\section}{0pt}{1.0em}{0.2em}
\titleformat{\subsection}{\color{navy}\bfseries\normalsize}{\thesubsection}{0.6em}{}
\titlespacing*{\subsection}{0pt}{0.9em}{0.15em}
% Indice: sezioni in blu scuro, sottosezioni con puntini
\titlecontents{section}[1.6em]{\addvspace{0.4em}\color{navy}\bfseries}{\contentslabel{1.6em}}{\hspace*{-1.6em}}{\hfill\contentspage}
\titlecontents{subsection}[4.0em]{}{\contentslabel{2.4em}}{\hspace*{-2.4em}}{\titlerule*[0.7pc]{.}\contentspage}
% Evita titoli orfani in fondo pagina: se restano poche righe si va a capo pagina
\pretocmd{\section}{\needspace{9\baselineskip}}{}{}
\pretocmd{\subsection}{\needspace{5\baselineskip}}{}{}
% Le tabelle restano nella sezione in cui sono citate: placeins[section] chiude i
% float a ogni \section; a fine sezione si aggiunge comunque \FloatBarrier esplicito.
\renewcommand{\topfraction}{0.9}
\renewcommand{\bottomfraction}{0.8}
\renewcommand{\textfraction}{0.07}
\renewcommand{\floatpagefraction}{0.75}
\setcounter{topnumber}{3}
\setcounter{bottomnumber}{2}
\setcounter{totalnumber}{5}

\setlist{itemsep=0.15em,topsep=0.2em,parsep=0pt,leftmargin=1.3em}

% Didascalie delle tabelle: sopra la tabella, "Tabella N." in blu grassetto
\DeclareCaptionFont{navy}{\color{navy}}
\captionsetup{font=small,labelfont={bf,navy},textfont=navy,labelsep=period,justification=raggedright,singlelinecheck=false,skip=4pt,position=top}
\setlength{\LTcapwidth}{\linewidth}

% Tabelle: intestazione blu con testo bianco, righe alternate grigie, filetti sottili.
% Uso: \begin{tabella} ... tabularx o xltabular ... \end{tabella};
% riga di intestazione: \rowcolor{navy}\hd{Colonna} & ... \\ \hline
% xltabular va a capo pagina e ripete l'intestazione posta prima di \endhead.
\arrayrulecolor{rulegrey}
\setlength{\arrayrulewidth}{0.4pt}
\setlength{\LTpre}{0.3em}
\setlength{\LTpost}{0.3em}
\newcommand{\hd}[1]{\textcolor{white}{\bfseries #1}}
\newcolumntype{L}{>{\raggedright\arraybackslash}X}
\newcolumntype{P}[1]{>{\raggedright\arraybackslash}p{#1}}
\newcolumntype{R}[1]{>{\raggedleft\arraybackslash}p{#1}}
\newcolumntype{C}[1]{>{\centering\arraybackslash}p{#1}}
\newenvironment{tabella}{\par\footnotesize\renewcommand{\arraystretch}{1.3}\rowcolors{2}{rowgrey}{white}}{\par\normalsize}

% Riquadro di sintesi: bordo blu, fondo grigio chiaro (non si spezza tra pagine)
\newsavebox{\riquadrobox}
\newenvironment{riquadro}
  {\par\noindent\begin{lrbox}{\riquadrobox}\begin{minipage}{\dimexpr\linewidth-2\fboxsep-2\fboxrule\relax}\vspace{0.3em}}
  {\vspace{0.2em}\end{minipage}\end{lrbox}\fcolorbox{navy}{boxgrey}{\usebox{\riquadrobox}}\par}

% Collegamenti: hyperref per ultimo
\usepackage{hyperref}
\hypersetup{colorlinks=true,linkcolor=navy,citecolor=navy,urlcolor=maroon}
\AtBeginDocument{\hypersetup{pdftitle={\ReportTitle},pdfauthor={\ReportAuthor}}}
% cleveref dopo hyperref: \cref{tab:x} produce "Tabella 1"
\usepackage[capitalise,noabbrev]{cleveref}
 e \begin{document}.
% Compilare con xelatex, due passaggi (tre se cambiano riferimenti o bibliografia).
\documentclass[10pt,a4paper]{article}

% Valori di default: il file principale li sovrascrive definendoli prima di \input.
\providecommand{\ReportLang}{italian}% italian | english
\providecommand{\ReportBrand}{Nomisma}
\providecommand{\ReportProject}{}
\providecommand{\ReportTitle}{<<TITOLO>>}
\providecommand{\ReportSubtitle}{<<SOTTOTITOLO>>}
\providecommand{\ReportAuthor}{<<AUTORE>>}
\providecommand{\ReportRole}{<<RUOLO>>}
\providecommand{\ReportDate}{<<DATA>>}

\usepackage{etoolbox}
\usepackage{fontspec}
% Testo estraibile corretto dal PDF: senza questa riga le parentesi di Segoe UI
% vengono copiate come caratteri ornamentali (U+FD3E/U+FD3F).
\XeTeXgenerateactualtext=1

% Font: Segoe UI di Windows come nella nota di riferimento; fuori da Windows si
% ripiega su Latin Modern Sans, sempre presente in MiKTeX e TeX Live.
\IfFileExists{C:/Windows/Fonts/segoeui.ttf}{%
  \setmainfont{segoeui}[
    Path = C:/Windows/Fonts/,
    Extension = .ttf,
    BoldFont = segoeuib,
    ItalicFont = segoeuii,
    BoldItalicFont = segoeuiz]
  \setsansfont{segoeui}[
    Path = C:/Windows/Fonts/,
    Extension = .ttf,
    BoldFont = segoeuib,
    ItalicFont = segoeuii,
    BoldItalicFont = segoeuiz]
}{%
  \setmainfont{Latin Modern Sans}
  \setsansfont{Latin Modern Sans}
}

\usepackage[english,italian]{babel}
\AtBeginDocument{\selectlanguage{\ReportLang}}

\usepackage[a4paper,top=2.3cm,bottom=2.2cm,left=2.2cm,right=2.2cm,headheight=14pt,headsep=0.6cm,footskip=1cm]{geometry}
\usepackage[table]{xcolor}
\usepackage{colortbl}
\usepackage{array}
\usepackage{tabularx}
\usepackage{xltabular}
\usepackage{titlesec}
\usepackage{fancyhdr}
\usepackage{enumitem}
\usepackage{microtype}
\usepackage{needspace}
\usepackage{titletoc}
\usepackage[section]{placeins}
\usepackage{amsmath}
\usepackage{amssymb}
\usepackage{graphicx}
\usepackage{caption}
\usepackage{siunitx}

% Colori della nota metodologica
\definecolor{navy}{HTML}{0E2D53}
\definecolor{maroon}{HTML}{A3274B}
\definecolor{rowgrey}{HTML}{F2F2F2}
\definecolor{rulegrey}{HTML}{BFBFBF}
\definecolor{footgrey}{HTML}{8C8C8C}
\definecolor{boxgrey}{HTML}{F4F5F8}

\setlength{\parindent}{0pt}
\setlength{\parskip}{0.55em}
\linespread{1.12}

% Numeri: in italiano punto per le migliaia e virgola per i decimali, in inglese
% il contrario. Nel sorgente si scrive sempre il numero grezzo: \num{161250},
% \num{0.698}, \pct{96.6}. Il raggruppamento parte da 4 cifre (1.626).
% Modo testo: le cifre usano il font del testo e ereditano il grassetto.
% Incertezza esplicita: \num{0.874 +- 0.006} produce 0,874 ± 0,006.
\sisetup{mode=text,reset-text-family=false,reset-text-series=false,reset-text-shape=false,uncertainty-mode=separate}
\ifdefstring{\ReportLang}{english}{%
  \sisetup{output-decimal-marker={.}, group-separator={,}, group-minimum-digits=4}
  \newcommand{\ReportPageLabel}{p.}
}{%
  \sisetup{output-decimal-marker={,}, group-separator={.}, group-minimum-digits=4}
  \newcommand{\ReportPageLabel}{pag.}
}
\newcommand{\pct}[1]{\num{#1}\%}

% Testatina bordeaux in maiuscolo a sinistra, piede grigio a destra
\providecommand{\ReportHeader}{\ReportBrand\ifdefempty{\ReportProject}{}{\ \ | \ \ReportProject}}
\providecommand{\ReportFooter}{\ReportBrand\ifdefempty{\ReportProject}{}{ | \ReportProject}}
\newcommand{\ReportHeadStyle}{\footnotesize\bfseries\color{maroon}\MakeUppercase{\ReportHeader}}
\newcommand{\ReportFootStyle}{\scriptsize\color{footgrey}\ReportFooter\ ·\ \ReportPageLabel\ \thepage}
\pagestyle{fancy}
\fancyhf{}
\renewcommand{\headrulewidth}{0pt}
\fancyhead[L]{\ReportHeadStyle}
\fancyfoot[R]{\ReportFootStyle}
\fancypagestyle{plain}{\fancyhf{}\renewcommand{\headrulewidth}{0pt}%
  \fancyhead[L]{\ReportHeadStyle}\fancyfoot[R]{\ReportFootStyle}}

% Blocco del titolo: titolo blu in grassetto, sottotitolo, autore e ruolo, data
\newcommand{\maketitleblock}{%
  {\color{navy}\bfseries\LARGE \ReportTitle\par}
  \vspace{0.35em}
  {\color{navy}\large \ReportSubtitle\par}
  \vspace{0.5em}
  {\small \ReportAuthor, \ReportRole\par
  \color{footgrey}\ReportDate\par}}

% Titoli di sezione: numero e titolo blu in grassetto con filetto bordeaux sotto.
% Le sottosezioni sono blu, in grassetto, di corpo normale (piu' piccole delle sezioni).
\newcommand{\sectitle}[1]{\parbox[b]{\linewidth}{\thesection.\hspace{0.4em}#1\par\nointerlineskip\vspace{3pt}{\color{maroon}\rule{\linewidth}{0.7pt}}}}
\newcommand{\sectitlenn}[1]{\parbox[b]{\linewidth}{#1\par\nointerlineskip\vspace{3pt}{\color{maroon}\rule{\linewidth}{0.7pt}}}}
\titleformat{\section}[block]{\color{navy}\bfseries\large}{}{0pt}{\sectitle}
\titleformat{name=\section,numberless}[block]{\color{navy}\bfseries\large}{}{0pt}{\sectitlenn}
\titlespacing*{\section}{0pt}{1.0em}{0.2em}
\titleformat{\subsection}{\color{navy}\bfseries\normalsize}{\thesubsection}{0.6em}{}
\titlespacing*{\subsection}{0pt}{0.9em}{0.15em}
% Indice: sezioni in blu scuro, sottosezioni con puntini
\titlecontents{section}[1.6em]{\addvspace{0.4em}\color{navy}\bfseries}{\contentslabel{1.6em}}{\hspace*{-1.6em}}{\hfill\contentspage}
\titlecontents{subsection}[4.0em]{}{\contentslabel{2.4em}}{\hspace*{-2.4em}}{\titlerule*[0.7pc]{.}\contentspage}
% Evita titoli orfani in fondo pagina: se restano poche righe si va a capo pagina
\pretocmd{\section}{\needspace{9\baselineskip}}{}{}
\pretocmd{\subsection}{\needspace{5\baselineskip}}{}{}
% Le tabelle restano nella sezione in cui sono citate: placeins[section] chiude i
% float a ogni \section; a fine sezione si aggiunge comunque \FloatBarrier esplicito.
\renewcommand{\topfraction}{0.9}
\renewcommand{\bottomfraction}{0.8}
\renewcommand{\textfraction}{0.07}
\renewcommand{\floatpagefraction}{0.75}
\setcounter{topnumber}{3}
\setcounter{bottomnumber}{2}
\setcounter{totalnumber}{5}

\setlist{itemsep=0.15em,topsep=0.2em,parsep=0pt,leftmargin=1.3em}

% Didascalie delle tabelle: sopra la tabella, "Tabella N." in blu grassetto
\DeclareCaptionFont{navy}{\color{navy}}
\captionsetup{font=small,labelfont={bf,navy},textfont=navy,labelsep=period,justification=raggedright,singlelinecheck=false,skip=4pt,position=top}
\setlength{\LTcapwidth}{\linewidth}

% Tabelle: intestazione blu con testo bianco, righe alternate grigie, filetti sottili.
% Uso: \begin{tabella} ... tabularx o xltabular ... \end{tabella};
% riga di intestazione: \rowcolor{navy}\hd{Colonna} & ... \\ \hline
% xltabular va a capo pagina e ripete l'intestazione posta prima di \endhead.
\arrayrulecolor{rulegrey}
\setlength{\arrayrulewidth}{0.4pt}
\setlength{\LTpre}{0.3em}
\setlength{\LTpost}{0.3em}
\newcommand{\hd}[1]{\textcolor{white}{\bfseries #1}}
\newcolumntype{L}{>{\raggedright\arraybackslash}X}
\newcolumntype{P}[1]{>{\raggedright\arraybackslash}p{#1}}
\newcolumntype{R}[1]{>{\raggedleft\arraybackslash}p{#1}}
\newcolumntype{C}[1]{>{\centering\arraybackslash}p{#1}}
\newenvironment{tabella}{\par\footnotesize\renewcommand{\arraystretch}{1.3}\rowcolors{2}{rowgrey}{white}}{\par\normalsize}

% Riquadro di sintesi: bordo blu, fondo grigio chiaro (non si spezza tra pagine)
\newsavebox{\riquadrobox}
\newenvironment{riquadro}
  {\par\noindent\begin{lrbox}{\riquadrobox}\begin{minipage}{\dimexpr\linewidth-2\fboxsep-2\fboxrule\relax}\vspace{0.3em}}
  {\vspace{0.2em}\end{minipage}\end{lrbox}\fcolorbox{navy}{boxgrey}{\usebox{\riquadrobox}}\par}

% Collegamenti: hyperref per ultimo
\usepackage{hyperref}
\hypersetup{colorlinks=true,linkcolor=navy,citecolor=navy,urlcolor=maroon}
\AtBeginDocument{\hypersetup{pdftitle={\ReportTitle},pdfauthor={\ReportAuthor}}}
% cleveref dopo hyperref: \cref{tab:x} produce "Tabella 1"
\usepackage[capitalise,noabbrev]{cleveref}
 e \begin{document}.
% Compilare con xelatex, due passaggi (tre se cambiano riferimenti o bibliografia).
\documentclass[10pt,a4paper]{article}

% Valori di default: il file principale li sovrascrive definendoli prima di \input.
\providecommand{\ReportLang}{italian}% italian | english
\providecommand{\ReportBrand}{Nomisma}
\providecommand{\ReportProject}{}
\providecommand{\ReportTitle}{<<TITOLO>>}
\providecommand{\ReportSubtitle}{<<SOTTOTITOLO>>}
\providecommand{\ReportAuthor}{<<AUTORE>>}
\providecommand{\ReportRole}{<<RUOLO>>}
\providecommand{\ReportDate}{<<DATA>>}

\usepackage{etoolbox}
\usepackage{fontspec}
% Testo estraibile corretto dal PDF: senza questa riga le parentesi di Segoe UI
% vengono copiate come caratteri ornamentali (U+FD3E/U+FD3F).
\XeTeXgenerateactualtext=1

% Font: Segoe UI di Windows come nella nota di riferimento; fuori da Windows si
% ripiega su Latin Modern Sans, sempre presente in MiKTeX e TeX Live.
\IfFileExists{C:/Windows/Fonts/segoeui.ttf}{%
  \setmainfont{segoeui}[
    Path = C:/Windows/Fonts/,
    Extension = .ttf,
    BoldFont = segoeuib,
    ItalicFont = segoeuii,
    BoldItalicFont = segoeuiz]
  \setsansfont{segoeui}[
    Path = C:/Windows/Fonts/,
    Extension = .ttf,
    BoldFont = segoeuib,
    ItalicFont = segoeuii,
    BoldItalicFont = segoeuiz]
}{%
  \setmainfont{Latin Modern Sans}
  \setsansfont{Latin Modern Sans}
}

\usepackage[english,italian]{babel}
\AtBeginDocument{\selectlanguage{\ReportLang}}

\usepackage[a4paper,top=2.3cm,bottom=2.2cm,left=2.2cm,right=2.2cm,headheight=14pt,headsep=0.6cm,footskip=1cm]{geometry}
\usepackage[table]{xcolor}
\usepackage{colortbl}
\usepackage{array}
\usepackage{tabularx}
\usepackage{xltabular}
\usepackage{titlesec}
\usepackage{fancyhdr}
\usepackage{enumitem}
\usepackage{microtype}
\usepackage{needspace}
\usepackage{titletoc}
\usepackage[section]{placeins}
\usepackage{amsmath}
\usepackage{amssymb}
\usepackage{graphicx}
\usepackage{caption}
\usepackage{siunitx}

% Colori della nota metodologica
\definecolor{navy}{HTML}{0E2D53}
\definecolor{maroon}{HTML}{A3274B}
\definecolor{rowgrey}{HTML}{F2F2F2}
\definecolor{rulegrey}{HTML}{BFBFBF}
\definecolor{footgrey}{HTML}{8C8C8C}
\definecolor{boxgrey}{HTML}{F4F5F8}

\setlength{\parindent}{0pt}
\setlength{\parskip}{0.55em}
\linespread{1.12}

% Numeri: in italiano punto per le migliaia e virgola per i decimali, in inglese
% il contrario. Nel sorgente si scrive sempre il numero grezzo: \num{161250},
% \num{0.698}, \pct{96.6}. Il raggruppamento parte da 4 cifre (1.626).
% Modo testo: le cifre usano il font del testo e ereditano il grassetto.
% Incertezza esplicita: \num{0.874 +- 0.006} produce 0,874 ± 0,006.
\sisetup{mode=text,reset-text-family=false,reset-text-series=false,reset-text-shape=false,uncertainty-mode=separate}
\ifdefstring{\ReportLang}{english}{%
  \sisetup{output-decimal-marker={.}, group-separator={,}, group-minimum-digits=4}
  \newcommand{\ReportPageLabel}{p.}
}{%
  \sisetup{output-decimal-marker={,}, group-separator={.}, group-minimum-digits=4}
  \newcommand{\ReportPageLabel}{pag.}
}
\newcommand{\pct}[1]{\num{#1}\%}

% Testatina bordeaux in maiuscolo a sinistra, piede grigio a destra
\providecommand{\ReportHeader}{\ReportBrand\ifdefempty{\ReportProject}{}{\ \ | \ \ReportProject}}
\providecommand{\ReportFooter}{\ReportBrand\ifdefempty{\ReportProject}{}{ | \ReportProject}}
\newcommand{\ReportHeadStyle}{\footnotesize\bfseries\color{maroon}\MakeUppercase{\ReportHeader}}
\newcommand{\ReportFootStyle}{\scriptsize\color{footgrey}\ReportFooter\ ·\ \ReportPageLabel\ \thepage}
\pagestyle{fancy}
\fancyhf{}
\renewcommand{\headrulewidth}{0pt}
\fancyhead[L]{\ReportHeadStyle}
\fancyfoot[R]{\ReportFootStyle}
\fancypagestyle{plain}{\fancyhf{}\renewcommand{\headrulewidth}{0pt}%
  \fancyhead[L]{\ReportHeadStyle}\fancyfoot[R]{\ReportFootStyle}}

% Blocco del titolo: titolo blu in grassetto, sottotitolo, autore e ruolo, data
\newcommand{\maketitleblock}{%
  {\color{navy}\bfseries\LARGE \ReportTitle\par}
  \vspace{0.35em}
  {\color{navy}\large \ReportSubtitle\par}
  \vspace{0.5em}
  {\small \ReportAuthor, \ReportRole\par
  \color{footgrey}\ReportDate\par}}

% Titoli di sezione: numero e titolo blu in grassetto con filetto bordeaux sotto.
% Le sottosezioni sono blu, in grassetto, di corpo normale (piu' piccole delle sezioni).
\newcommand{\sectitle}[1]{\parbox[b]{\linewidth}{\thesection.\hspace{0.4em}#1\par\nointerlineskip\vspace{3pt}{\color{maroon}\rule{\linewidth}{0.7pt}}}}
\newcommand{\sectitlenn}[1]{\parbox[b]{\linewidth}{#1\par\nointerlineskip\vspace{3pt}{\color{maroon}\rule{\linewidth}{0.7pt}}}}
\titleformat{\section}[block]{\color{navy}\bfseries\large}{}{0pt}{\sectitle}
\titleformat{name=\section,numberless}[block]{\color{navy}\bfseries\large}{}{0pt}{\sectitlenn}
\titlespacing*{\section}{0pt}{1.0em}{0.2em}
\titleformat{\subsection}{\color{navy}\bfseries\normalsize}{\thesubsection}{0.6em}{}
\titlespacing*{\subsection}{0pt}{0.9em}{0.15em}
% Indice: sezioni in blu scuro, sottosezioni con puntini
\titlecontents{section}[1.6em]{\addvspace{0.4em}\color{navy}\bfseries}{\contentslabel{1.6em}}{\hspace*{-1.6em}}{\hfill\contentspage}
\titlecontents{subsection}[4.0em]{}{\contentslabel{2.4em}}{\hspace*{-2.4em}}{\titlerule*[0.7pc]{.}\contentspage}
% Evita titoli orfani in fondo pagina: se restano poche righe si va a capo pagina
\pretocmd{\section}{\needspace{9\baselineskip}}{}{}
\pretocmd{\subsection}{\needspace{5\baselineskip}}{}{}
% Le tabelle restano nella sezione in cui sono citate: placeins[section] chiude i
% float a ogni \section; a fine sezione si aggiunge comunque \FloatBarrier esplicito.
\renewcommand{\topfraction}{0.9}
\renewcommand{\bottomfraction}{0.8}
\renewcommand{\textfraction}{0.07}
\renewcommand{\floatpagefraction}{0.75}
\setcounter{topnumber}{3}
\setcounter{bottomnumber}{2}
\setcounter{totalnumber}{5}

\setlist{itemsep=0.15em,topsep=0.2em,parsep=0pt,leftmargin=1.3em}

% Didascalie delle tabelle: sopra la tabella, "Tabella N." in blu grassetto
\DeclareCaptionFont{navy}{\color{navy}}
\captionsetup{font=small,labelfont={bf,navy},textfont=navy,labelsep=period,justification=raggedright,singlelinecheck=false,skip=4pt,position=top}
\setlength{\LTcapwidth}{\linewidth}

% Tabelle: intestazione blu con testo bianco, righe alternate grigie, filetti sottili.
% Uso: \begin{tabella} ... tabularx o xltabular ... \end{tabella};
% riga di intestazione: \rowcolor{navy}\hd{Colonna} & ... \\ \hline
% xltabular va a capo pagina e ripete l'intestazione posta prima di \endhead.
\arrayrulecolor{rulegrey}
\setlength{\arrayrulewidth}{0.4pt}
\setlength{\LTpre}{0.3em}
\setlength{\LTpost}{0.3em}
\newcommand{\hd}[1]{\textcolor{white}{\bfseries #1}}
\newcolumntype{L}{>{\raggedright\arraybackslash}X}
\newcolumntype{P}[1]{>{\raggedright\arraybackslash}p{#1}}
\newcolumntype{R}[1]{>{\raggedleft\arraybackslash}p{#1}}
\newcolumntype{C}[1]{>{\centering\arraybackslash}p{#1}}
\newenvironment{tabella}{\par\footnotesize\renewcommand{\arraystretch}{1.3}\rowcolors{2}{rowgrey}{white}}{\par\normalsize}

% Riquadro di sintesi: bordo blu, fondo grigio chiaro (non si spezza tra pagine)
\newsavebox{\riquadrobox}
\newenvironment{riquadro}
  {\par\noindent\begin{lrbox}{\riquadrobox}\begin{minipage}{\dimexpr\linewidth-2\fboxsep-2\fboxrule\relax}\vspace{0.3em}}
  {\vspace{0.2em}\end{minipage}\end{lrbox}\fcolorbox{navy}{boxgrey}{\usebox{\riquadrobox}}\par}

% Collegamenti: hyperref per ultimo
\usepackage{hyperref}
\hypersetup{colorlinks=true,linkcolor=navy,citecolor=navy,urlcolor=maroon}
\AtBeginDocument{\hypersetup{pdftitle={\ReportTitle},pdfauthor={\ReportAuthor}}}
% cleveref dopo hyperref: \cref{tab:x} produce "Tabella 1"
\usepackage[capitalise,noabbrev]{cleveref}


\begin{document}

\maketitleblock

\section{In sintesi}

% GUIDA: due o tre frasi: che cosa è stato fatto, su quale scala, con quale sforzo (volumi, fonti incrociate, verifiche).
<<Che cosa abbiamo fatto, per quante unità, incrociando quante fonti e con quante verifiche.>>

% GUIDA: da 3 a 5 punti con i numeri principali; nessun codice, nessun nome di variabile.
\begin{riquadro}
\begin{itemize}
\item \textbf{<<Etichetta>>}: <<numero principale e che cosa significa>>.
\item \textbf{<<Etichetta>>}: <<secondo risultato>>.
\item \textbf{<<Etichetta>>}: <<verifiche svolte e loro esito>>.
\end{itemize}
\end{riquadro}

% GUIDA: l'indice va sempre subito dopo il riquadro di sintesi.
\tableofcontents

\section{Il punto di partenza}

% GUIDA: che cosa: la domanda del cliente, l'oggetto, la definizione adottata e il vincolo principale, in parole semplici.
<<La richiesta, la definizione adottata e che cosa doveva garantire il risultato.>>

\section{Perché conta}

% GUIDA: perché: la rilevanza per il lettore e la decisione che il risultato sostiene. Se è breve, si può unire alla sezione precedente.
<<A chi serve il risultato e quale decisione sostiene.>>

\section{Come funziona il metodo}

% GUIDA: come: è la parte più completa del rapporto. Una sottosezione per blocco; ogni idea tecnica spiegata con parole comuni.
<<Il metodo in una frase.>>

\subsection{Le fonti dei dati}

% GUIDA: perché nessuna fonte basta da sola; poi la tabella delle fonti, citata nel testo con Tabella~\ref.
<<Perché servono più fonti e come sono state scelte; la Tabella~\ref{tab:fonti} riporta per ciascuna che cosa contiene, perché la usiamo, data e copertura.>>

% GUIDA: quali alternative sono state valutate e scartate, e perché (prova dell'ampiezza della ricerca).
<<Alternative considerate e scartate, con il motivo.>>

% GUIDA: tabella lunga: xltabular con didascalia e intestazione ripetuta; la tabella breve va invece in un float (vedi Risultati).
\begin{tabella}
\begin{xltabular}{\linewidth}{|P{3.0cm}|L|P{3.3cm}|P{4.0cm}|}
\rowcolor{white}\caption{<<Le fonti dei dati>>}\label{tab:fonti}\\
\hline
\rowcolor{navy}\hd{Fonte} & \hd{Che cosa contiene} & \hd{Perché la usiamo} & \hd{Copertura e note}\\ \hline
\endfirsthead
\hline
\rowcolor{navy}\hd{Fonte} & \hd{Che cosa contiene} & \hd{Perché la usiamo} & \hd{Copertura e note}\\ \hline
\endhead
<<Fonte>> & <<Contenuto>> & <<Criterio che alimenta>> & <<Licenza; aggiornamento; unità abbinate>>\\ \hline
\end{xltabular}
\end{tabella}

\subsection{I passi, nell'ordine}

% GUIDA: passi numerati in linguaggio quotidiano, ciascuno con il suo perché; i codici delle regole compaiono solo in tabella.
\textbf{1. <<Primo passo>>.} <<Che cosa fa e perché.>>

\textbf{2. <<Secondo passo>>.} <<Che cosa fa e perché.>>

\subsection{Le soglie e la loro taratura}

% GUIDA: ogni soglia con il suo valore, il motivo, le alternative considerate e che cosa avrebbero prodotto; più esplicita dove la scelta sembra strana o asimmetrica.
<<Valore della soglia, perché è stata scelta, alternative e loro effetto.>>

\subsection{I controlli}

% GUIDA: taratura, controlli incrociati, revisione dei casi dubbi; ogni indice spiegato a parole (vedi esempio del kappa).
<<Come è stato controllato il risultato e con quale esito.>>

% GUIDA: una frase onesta sul metodo se sono stati usati strumenti assistiti dall'intelligenza artificiale.
<<Le verifiche mirate sono state svolte con strumenti di ricerca assistiti dall'intelligenza artificiale; ogni giudizio è registrato con la sua fonte.>>

\FloatBarrier

\section{Risultati}

% GUIDA: prima la tabella, poi la lettura in due o tre frasi; percentuali da \pct, conteggi da \num. Ogni tabella sta in un float con \caption e \label ed è citata con Tabella~\ref.
La Tabella~\ref{tab:risultati} riporta <<che cosa mostra la tabella>>.

\begin{table}[!htbp]
\begin{tabella}
\caption{<<Che cosa mostra la tabella dei risultati>>}\label{tab:risultati}
\begin{tabularx}{\linewidth}{|L|R{2.0cm}|R{1.5cm}|}
\hline
\rowcolor{navy}\hd{<<Categoria>>} & \hd{<<Unità>>} & \hd{\%}\\ \hline
<<Voce>> & \num{0} & \num{0.0}\\ \hline
\rowcolor{white}\textbf{Totale} & \textbf{\num{0}} & \textbf{\num{100.0}}\\ \hline
\end{tabularx}
\end{tabella}
\end{table}

<<Che cosa dicono i numeri, con il confronto più informativo.>>

\FloatBarrier

\section{Limiti e come leggere i risultati}

% GUIDA: ogni limite con la sua dimensione (quante unità, quale quota) e come tenerne conto.
\begin{itemize}
\item \textbf{<<Limite>>.} <<Quanto pesa e come leggerlo.>>
\end{itemize}

\FloatBarrier

\section{Materiali consegnati}

% GUIDA: elenco dei file consegnati; il codice si cita di passaggio e rinvia al suo README.
\begin{itemize}
\item \textbf{<<Il file principale>>} \texttt{<<nome\_file.xlsx>>}: <<che cosa contiene>>.
\item \textbf{<<Il pacchetto di codice>>} \texttt{<<nome\_pacchetto.zip>>}: <<permette di riprodurre il processo>>; installazione e uso sono descritti nel README contenuto nell'archivio.
\end{itemize}

\FloatBarrier

% GUIDA: solo per edizioni riviste, sempre ultima sezione. Ogni confronto con la versione precedente, vecchia soglia o vecchio numero sta qui e solo qui; il resto del rapporto si legge come se non esistesse una versione precedente. Per un'edizione unica, cancellare il blocco.
% \section{Note di versione}
%
% <<Che cosa è cambiato rispetto alla versione precedente e perché.>>
%
% \subsection{<<Il confronto dei risultati>>}
%
% <<Confronto con i numeri della versione precedente.>>
%
% \FloatBarrier

\end{document}
